\begin{abstract}
Composed video retrieval combines a reference video and modification text to retrieve a target, yet conventional evaluation can conflate edit-direction understanding with shortcuts from visual similarity, speech transcripts, or easy candidate pools. We introduce \textbf{Audio-CVR}, which studies directional sound modifications under preserved visual context: the reference must not satisfy the edit, the target must satisfy it, and muted video alone must not identify the target. We develop a multi-stage multimodal curation pipeline comprising source-aware clip pairing, audio-only difference analysis and edit generation, directional verification, muted-video shortcut rejection, full audiovisual consistency checking, anti-ASR filtering, repeat auditing, and source-disjoint splitting. We further establish a reference-aware protocol that treats the unchanged reference as a mandatory counterfactual negative and combines it with typed hard negatives and random distractors. Controlled reference and seven-mode modality ablations measure recall, target-over-reference accuracy, directional score margins, and reference-induced R@1 drop. We evaluate a frozen E5-Omni-7B backbone with a lightweight low-rank adapter on an automatically curated and model-verified non-speech benchmark. Across 150 source-disjoint queries, fixed 1,000-video galleries, and five seeds, enabling audio raises R@1 from 11.3\% to 22.9\%, a paired gain of 11.6 points (95\% CI: 6.3--17.1; Holm-adjusted $p<0.001$), and improves target-over-reference accuracy from 11.3\% to 23.2\%. Removing the unchanged reference raises R@1 to 99.5\%, showing that conventional galleries can conceal the difficulty of edit-direction discrimination. Audio-CVR therefore makes directional audio use, rather than generic video similarity, an explicit and measurable retrieval requirement.
\end{abstract}
